\section{Challenges and Future Directions}\label{sec:dis}
\subsection{Current Challenges}\label{sec:challenges}
Graph prompt learning has made significant research progress, but it still encounters several challenges. In this subsection, we will discuss current challenges in detail.

\textbf{Inherent Limitation within Graph Models} Prompts, originated from the NLP field \cite{brown2020language, lester2021power, liu2023pretrain}, serve as a means to unlock the potential of pre-trained language models for adapting to downstream tasks. This in-context learning ability emerges when the scale of model parameters reaches a certain level. For instance, widely known LLMs typically possess billions of parameters. With such a powerful pre-trained model, a simple textual prompt can distill specific knowledge for downstream tasks. However, when it comes to prompting on graph tasks, the conditions become more intractable since pre-trained graph models have significantly fewer parameters, making it challenging to harness their full downstream adaptation potential. 


\textbf{Intuitive Evaluation of Graph Prompt Learning}
In the NLP field, prompts are typically in a discrete textual format \cite{brown2020language, robinson2023leveraging}, allowing for intuitive understanding, comparison, and explanation. However, existing graph prompts are represented as learnable tokens \cite{sun2022gppt, fang2023universal, fang2022prompt, tan2023virtual, ma2023hetgpt} or augmented graphs \cite{sun2023all, huang2023prodigy, ge2023enhancing}. Such format poses challenges in intuitively understanding and interpreting graph prompts, as they lack a readable design. As a result, the effectiveness of prompts can only be evaluated based on downstream tasks, limiting efficient and comprehensive performance comparison among different kinds of graph prompts. Therefore, the development of a more intuitive graph prompt design with a readable format remains an open problem.


\textbf{More Downstream Applications}
Currently, graph prompt learning is primarily applied to node or graph classification tasks on open-source benchmarks \cite{hamilton2017inductive, xu2018how}. Although potential applications have been discussed in Section \ref{sec:app}, the real-world utilization of graph prompt learning remains limited. A notable example is its use in fraud detection within real-world transaction networks, addressing the issue of label scarcity \cite{wen2023voucher}. However, compared to the widespread application of prompting techniques in the NLP domain \cite{bran2023transformers,  robinson2023leveraging, zhang2023benchmarking}, the potential of graph prompt learning in diverse real-world applications requires further exploration.  Overcoming the main challenges of obtaining powerful domain-specific pre-trained graph models and designing suitable prompts for specific application scenarios that exhibit unique characteristics remains crucial.



\textbf{Transferable Prompt Designs} 
Existing studies on graph prompt learning \cite{liu2023graphprompt, sun2022gppt, gong2023prompt, fang2023universal, fang2022prompt, ge2023enhancing} typically focus on pre-training and prompt tuning using the same dataset, which limits the exploration of more transferable designs and empirical evaluation. Although PRODIGY \cite{huang2023prodigy} explores transferability within the same domain and All in One \cite{sun2023all} provides empirical results regarding transferability across tasks and domains, the investigation of prompt learning across diverse domains and tasks remains limited. Achieving transferability across diverse tasks and domains requires aligning the task space \cite{sun2023all}, semantic features \cite{zhu2023graphcontrol}, and structural patterns \cite{zhao2023graphglow}, which necessitates further theoretical work to provide insights guiding the development of transferable prompt designs.




\subsection{Future Directions}
With the above analysis on graph prompt, we summarize future directions as follows:

\textbf{Learning Knowledge from Large Graph Models (LGMs) like LLMs.} Currently, graph models are typically tailored to specific domains or tasks, limiting generalization abilities across broader domains or tasks \cite{liu2023graph, zhu2023graphcontrol}. Therefore, we are expecting the realization of large graph models (LGMs) as a universal tool to be intelligent enough to handle graph tasks across domains \cite{liu2023one, fatemi2023talk}. A significant step towards this direction has been taken by \citet{liu2023one}, who proposed a powerful model (OFA) capable of addressing classification tasks for graphs coming from various domains. However, their approach still relies on LLMs as a general template to distill specific domain knowledge, which we believe can be replaced with suitable graph prompting techniques. Just as textual prompt has been widely used to adapt LLM for diverse applications \cite{robinson2023leveraging, wang2023knowledge, brown2020language}, graph prompting techniques are promising to distill LGMs' knowledge specific for concrete downstream tasks. Applying the graph prompts to LGMs has the potential to revolutionize the field of deep graph learning, leveraging LGMs as a universal tool to tackle different tasks on graphs from diverse domains. 


\textbf{Transferable Learning.} As we have discussed in Section \ref{sec:challenges}, current graph prompt learning methods \cite{sun2023all, sun2022gppt, huang2023prodigy, liu2023graphprompt, tan2023virtual, ge2023enhancing, shirkavand2023deep} are primarily limited to intra-dataset/domain settings, lacking the ability to transfer knowledge across different tasks or domains. While there have been efforts in graph transfer learning \cite{zhu2023graphcontrol} to realize domain adaptation, research specifically focused on transferable graph prompting techniques remains limited. To enable the transfer ability, the prompts should be designed to realize the alignment between different task spaces \cite{sun2023all}, reformulating different tasks on graphs into a uniform template. Besides, both structural \cite{zhao2023graphglow} and semantics alignment \cite{zhu2023graphcontrol} should be realized via suitable graph prompts to enable domain adaptation. It is promising to implement transferable graph prompts, realizing knowledge transfer across domains to extend the generalization and applicability of graph models.




\textbf{More Theoretical Foundation.} Despite the great success of graph prompt learning on various downstream tasks, they mostly draw on the successful experience of prompt tuning on NLP domain \cite{liu2022ptuning, brown2020language, gao2021making, tsimpoukelli2021multimodal, qin2021learning, shin2020autoprompt}. In other words, most existing graph prompt tuning methods are designed with intuition, and their performance gains are evaluated by empirical experiments. The lack of sufficient theoretical foundations behind the design has led to both performance bottlenecks and poor explainability. Therefore, we believe that building a solid theoretical foundation for graph prompt learning from a graph theory perspective and minimizing the gap between the theoretical foundation and empirical design is also a promising future direction.


\textbf{More Explainable and Understandable Design. }While existing graph prompt learning methods have demonstrated impressive results on various downstream tasks \cite{guo2023datacentric, wen2023voucher, yang2023empirical}, we still lack a clear understanding of what exactly is being learned from the prompts. The black-box learning mode of prompt vectors raises questions about their interpretability and whether we can establish meaningful correspondences between the input data and the prompted graph \cite{sun2023all}. These issues are crucial for understanding and interpreting prompts but are currently missing in most graph prompt research. To work towards trustworthy graph prompt learning, it is promising to explore the self-interpretability \cite{dai2021selfexplainable}
% , zhang2021protgnn} 
to enable intuitive explanations of graph prompts. By gaining insights into the learned prompt vectors and structures, we can enhance our understanding of the underlying mechanisms and improve the interpretability of graph prompts. This, in turn, can lead to more effective utilization of prompts for security- or privacy-related downstream tasks.